% ============================================================
\chapter{TDA et Apprentissage Profond}
\label{ch:deep_tda}
% ============================================================

\begin{intuition}
L'int\'egration de la TDA dans l'apprentissage profond pose un
d\'efi fondamental~: le calcul de la persistance n'est a priori
\textbf{pas diff\'erentiable}. Les travaux r\'ecents ont montr\'e
comment rendre la persistance diff\'erentiable, permettant
d'entra\^\i ner des r\'eseaux de neurones de bout en bout avec
des objectifs topologiques.
\end{intuition}

% ============================================================
\section{Persistance diff\'erentiable}
\label{sec:persistance_diff}
% ============================================================

\begin{definition}[Persistance comme fonction des filtrations]
\index{persistance diff\'erentiable}
Consid\'erons un complexe simplicial $K$ avec une filtration
param\'etrique $f_\theta : K \to \R$ d\'ependant de param\`etres
$\theta \in \R^p$. Le diagramme de persistance
$\mathrm{Dgm}(f_\theta)$ est une fonction de $\theta$.
\end{definition}

\begin{theoreme}[Diff\'erentiabilit\'e --- Gameiro et al.\ 2016, Poulenard et al.\ 2018]
\label{thm:diff_persistance}
\index{diff\'erentiabilit\'e de la persistance}
Soit $K$ un complexe simplicial fini et $f : K \to \R$ une
filtration g\'en\'erique (valeurs de filtration distinctes).
Les coordonn\'ees de naissance $b_i(\theta)$ et de mort
$d_i(\theta)$ dans le diagramme de persistance sont des fonctions
diff\'erentiables de $\theta$ presque partout. Plus pr\'ecis\'ement~:
\[
  \frac{\partial b_i}{\partial \theta_j} =
  \frac{\partial f_\theta(\sigma_{b_i})}{\partial \theta_j}, \qquad
  \frac{\partial d_i}{\partial \theta_j} =
  \frac{\partial f_\theta(\sigma_{d_i})}{\partial \theta_j}
\]
o\`u $\sigma_{b_i}$ et $\sigma_{d_i}$ sont les simplexes qui
cr\'eent et tuent la classe $i$.
\end{theoreme}

\begin{attention}[title=\textbf{Points de non-diff\'erentiabilit\'e}]
La diff\'erentiabilit\'e \'echoue quand deux valeurs de filtration
co\"\i ncident (les simplexes de naissance et de mort peuvent
changer). En pratique, on ajoute une petite perturbation ou on
utilise des sous-gradients.
\end{attention}

\begin{formules}[title=\textbf{Gradient de la persistance}]
\[
  \frac{\partial \mathrm{pers}_i}{\partial \theta}
  = \frac{\partial d_i}{\partial \theta}
  - \frac{\partial b_i}{\partial \theta}
  = \nabla_\theta f(\sigma_{d_i}) - \nabla_\theta f(\sigma_{b_i})
\]
\end{formules}

\begin{minted}{python}
import torch
from topologylayer.nn import RipsLayer

# Nuage de points differentiable
X = torch.randn(50, 2, requires_grad=True)

# Couche de persistance differentiable
rips_layer = RipsLayer(maxdim=1, sublevel=False)
dgm = rips_layer(X)

# Loss basee sur la persistance H1
loss = -dgm[1][:, 1].sum()  # maximiser les persistances H1
loss.backward()
print(f"Gradient shape: {X.grad.shape}")
\end{minted}

% ============================================================
\section{PersLay~: couche de persistance g\'en\'erique}
\label{sec:perslay}
% ============================================================

\begin{definition}[PersLay --- Carri\`ere et al.\ 2020]
\index{PersLay}
\textbf{PersLay} est une couche neuronale qui transforme un
diagramme de persistance $D = \{(b_i, d_i)\}_{i=1}^n$ en un
vecteur de dimension fixe via~:
\[
  \mathrm{PersLay}(D) = \mathrm{op}\!\left(
  w(p_i) \cdot \phi(p_i)\right)_{i=1}^n
\]
o\`u~:
\begin{itemize}
  \item $w : \R^2 \to \R_+$ est une fonction de poids
        (att\'enuation des points proches de la diagonale).
  \item $\phi : \R^2 \to \R^q$ est une fonction de
        repr\'esentation point par point (apprise ou fix\'ee).
  \item $\mathrm{op}$ est une op\'eration d'agr\'egation
        (somme, max, top-$k$, attention).
\end{itemize}
\end{definition}

\begin{proposition}[Universalit\'e de PersLay]
Avec des choix ad\'equats de $w$, $\phi$ et $\mathrm{op}$,
PersLay peut approcher uniform\'ement toute fonction continue
et invariante par permutation sur les diagrammes de persistance.
\end{proposition}

\begin{remarque}
PersLay unifie les repr\'esentations classiques~:
\begin{itemize}
  \item \textbf{Paysages}~: $\phi$ est la fonction tente,
        $\mathrm{op}$ est top-$k$.
  \item \textbf{Images}~: $\phi$ est une gaussienne,
        $\mathrm{op}$ est la somme.
  \item \textbf{Silhouettes}~: $\phi$ est la fonction tente,
        $\mathrm{op}$ est la moyenne pond\'er\'ee.
\end{itemize}
\end{remarque}

% ============================================================
\section{R\'egularisation topologique}
\label{sec:regularisation}
% ============================================================

\begin{definition}[Perte topologique]
\index{r\'egularisation topologique}
Une \textbf{perte topologique} est un terme ajout\'e \`a la
fonction de co\^ut pour encourager ou p\'enaliser certaines
structures topologiques. Formellement~:
\[
  \mathcal{L}_{\mathrm{total}} = \mathcal{L}_{\mathrm{task}}
  + \lambda \cdot \mathcal{L}_{\mathrm{topo}}
\]
o\`u $\mathcal{L}_{\mathrm{topo}}$ est calcul\'ee \`a partir
du diagramme de persistance.
\end{definition}

\begin{exemple}[Perte topologique pour la segmentation]
En segmentation d'images, on peut p\'enaliser les petits trous
parasites dans la pr\'ediction~:
\[
  \mathcal{L}_{\mathrm{topo}} = \sum_{(b_i, d_i) \in H_1}
  (d_i - b_i)^2
\]
qui force le r\'eseau \`a produire des r\'egions simplement connexes.
\end{exemple}

\begin{theoreme}[Convergence de la r\'egularisation topologique]
\label{thm:convergence_topo}
Sous des conditions de r\'egularit\'e (filtration g\'en\'erique,
Lipschitz), la descente de gradient stochastique avec perte
topologique converge vers un point critique de
$\mathcal{L}_{\mathrm{total}}$.
\end{theoreme}

\begin{minted}{python}
import torch
import torch.nn as nn

class TopologicalLoss(nn.Module):
    """Penalise les trous parasites en segmentation."""
    def __init__(self, lam=1.0):
        super().__init__()
        self.lam = lam
        # rips_layer defini avec une couche differentiable

    def forward(self, pred, target, dgm_h1):
        ce_loss = nn.functional.cross_entropy(pred, target)
        # Persistances H1 = d - b
        pers = dgm_h1[:, 1] - dgm_h1[:, 0]
        topo_loss = (pers ** 2).sum()
        return ce_loss + self.lam * topo_loss
\end{minted}

% ============================================================
\section{Auto-encodeurs topologiques}
\label{sec:autoencodeurs_topo}
% ============================================================

\begin{definition}[Auto-encodeur topologique --- Moor et al.\ 2020]
\index{auto-encodeur topologique}
Un \textbf{auto-encodeur topologique} est un auto-encodeur
$(\mathrm{Enc}, \mathrm{Dec})$ dont la perte inclut un terme
topologique~:
\[
  \mathcal{L} = \norm{x - \mathrm{Dec}(\mathrm{Enc}(x))}^2
  + \lambda \cdot d_W\!\left(\mathrm{Dgm}(X),\;
  \mathrm{Dgm}(\mathrm{Enc}(X))\right)^2
\]
o\`u $d_W$ est une distance de Wasserstein entre les diagrammes
de persistance de l'espace d'entr\'ee $X$ et de l'espace latent
$\mathrm{Enc}(X)$.
\end{definition}

\begin{intuition}
L'id\'ee est de pr\'eserver la \textbf{topologie} des donn\'ees
lors de la r\'eduction de dimension. Un auto-encodeur classique (ou
un VAE) peut \'ecraser des boucles ou fusionner des composantes
connexes. Le terme topologique p\'enalise ces distorsions.
\end{intuition}

\begin{proposition}[Pr\'eservation topologique]
Si $\mathcal{L}_{\mathrm{topo}} = 0$, alors les nombres de Betti
de $X$ et de $\mathrm{Enc}(X)$ co\"\i ncident pour la filtration
de Rips utilis\'ee.
\end{proposition}

% ============================================================
\section{R\'eseaux de neurones sur graphes et homologie persistante}
\label{sec:gnn_persistance}
% ============================================================

\begin{definition}[Filtration apprise sur un graphe]
\index{filtration apprise}
Soit $G = (V, E)$ un graphe avec des features de n\oe uds
$h_v \in \R^d$. On d\'efinit une filtration apprise~:
\[
  f_\theta(v) = g_\theta(h_v), \qquad
  f_\theta(e_{uv}) = \max(f_\theta(u), f_\theta(v))
\]
o\`u $g_\theta : \R^d \to \R$ est un r\'eseau de neurones.
Le diagramme de persistance de cette filtration fournit des
features topologiques du graphe.
\end{definition}

\begin{exemple}[Classification de graphes]
Pour la classification de graphes mol\'eculaires~:
\begin{enumerate}
  \item Un GNN produit des embeddings de n\oe uds $h_v$.
  \item Une filtration apprise $f_\theta$ est calcul\'ee.
  \item Le diagramme de persistance est vectoris\'e via PersLay.
  \item Le vecteur r\'esultant est concat\'en\'e au readout du GNN.
\end{enumerate}
Cette combinaison am\'eliore l'expressivit\'e au-del\`a du test
de Weisfeiler-Leman.
\end{exemple}

\begin{theoreme}[Expressivit\'e topologique --- Horn et al.\ 2022]
\index{expressivit\'e}
Il existe des paires de graphes non distinguables par le test de
Weisfeiler-Leman de niveau $k$ mais distinguables par la
persistance de la filtration induite par un GNN de message passing.
\end{theoreme}

% ============================================================
\section{Impl\'ementation avec PyTorch}
\label{sec:implementation}
% ============================================================

\begin{minted}{python}
import torch
import torch.nn as nn
from torch_geometric.nn import GCNConv, global_mean_pool

class TopoGNN(nn.Module):
    """GNN avec features topologiques."""
    def __init__(self, in_dim, hidden_dim, out_dim):
        super().__init__()
        self.conv1 = GCNConv(in_dim, hidden_dim)
        self.conv2 = GCNConv(hidden_dim, hidden_dim)
        self.filtration = nn.Linear(hidden_dim, 1)
        self.classifier = nn.Linear(hidden_dim + 16, out_dim)
        # perslay: couche de vectorisation apprise
        # self.perslay = PersLay(...)

    def forward(self, x, edge_index, batch):
        h = torch.relu(self.conv1(x, edge_index))
        h = torch.relu(self.conv2(h, edge_index))
        # Readout classique
        graph_emb = global_mean_pool(h, batch)
        # Filtration apprise
        filt_vals = self.filtration(h).squeeze()
        # ... calcul de persistance et PersLay ...
        # topo_features = self.perslay(dgm)
        # out = self.classifier(
        #     torch.cat([graph_emb, topo_features], dim=-1)
        # )
        return graph_emb  # simplifie
\end{minted}

% ============================================================
\section{Exercices}
% ============================================================

\begin{exercice}
Calculer \`a la main le gradient de la persistance totale
$\sum_i (d_i - b_i)$ par rapport aux valeurs de filtration pour
le complexe $\{a, b, c, ab, bc\}$ avec $f(a) = 0, f(b) = 1,
f(c) = 2, f(ab) = 1, f(bc) = 2$.
\end{exercice}

\begin{exercice}
Montrer que PersLay avec $\phi(b, d) = (b, d-b)$,
$w \equiv 1$ et $\mathrm{op} = \mathrm{sum}$ revient \`a
calculer $(\sum b_i, \sum(d_i - b_i))$.
\end{exercice}

\begin{exercice}[Impl\'ementation]
Impl\'ementer un auto-encodeur topologique simple sur un jeu de
donn\'ees en forme de cercle. Comparer l'espace latent avec et
sans le terme topologique.
\end{exercice}

\begin{exercice}
Donner un exemple de deux graphes non isomorphes indistinguables
par $1$-WL mais distinguables par la persistance $H_0$ d'une
filtration par degr\'e.
\end{exercice}

\begin{exercice}[Projet]
Entra\^\i ner un r\'eseau de segmentation d'image (U-Net) avec
et sans r\'egularisation topologique sur un jeu de donn\'ees de
vaisseaux sanguins. Comparer les nombres de Betti des
pr\'edictions.
\end{exercice}
